% Source: source/普通高考/2010/2010广东理.pdf, page 1, question 13.
% Complete flowchart transcription from the source PDF.
\begin{tikzpicture}[
  >=Stealth,
  x=1cm,
  y=.85cm,
  line width=.45pt,
  line cap=round,
  line join=round,
  font=\normalsize,
  flowtext/.style={anchor=center,align=center,
    execute at begin node={\setlength{\parindent}{0pt}\noindent}},
  every node/.style={flowtext},
  flowbox/.style={flowtext,draw,minimum height=.70cm,
    inner xsep=4pt,inner ysep=2.5pt},
  terminal/.style={flowbox,rounded corners=.18cm,minimum width=1.45cm},
  process/.style={flowbox},
  decision/.style={flowbox,diamond,aspect=2.8,minimum width=4.35cm,
    minimum height=.95cm,inner sep=1pt},
  io/.style={flowbox,trapezium,trapezium left angle=72,
    trapezium right angle=108,minimum height=.78cm},
  flowarrow/.style={->,line width=.45pt}
]
  \node[terminal] (start) at (0,11.75) {开始};
  \node[io,minimum width=4.20cm] (input) at (0,10.30)
    {输入 $n,x_1,x_2,\ldots,x_n$};
  \node[process,minimum width=4.20cm] (init) at (0,8.90)
    {$s_1=0,\ s_2=0,\ i=1$};
  \coordinate (loopjoin) at (0,7.90);

  \node[process,minimum width=2.45cm] (inc) at (3.95,7.90)
    {$i=i+1$};
  \node[process,minimum width=4.20cm,minimum height=1.05cm] (formula)
    at (3.95,6.05) {$s=\dfrac{1}{i}\left(s_2-\dfrac{1}{i}s_1^2\right)$};
  \node[process,minimum width=3.20cm,minimum height=1.10cm] (update)
    at (3.95,4.00) {$s_1=s_1+x_i$\\$s_2=s_2+x_i^2$};

  \node[decision] (test) at (0,4.00) {$i\leqslant n$};
  \node[io,minimum width=2.25cm] (output) at (0,2.35) {输出 $s$};
  \node[terminal] (finish) at (0,.95) {结束};

  \draw[flowarrow] (start.south) -- (input.north);
  \draw[flowarrow] (input.south) -- (init.north);
  \draw (init.south) -- (loopjoin);
  \draw[flowarrow] (loopjoin) -- (test.north);

  \draw[flowarrow] (test.east) -- node[above=2pt] {是} (update.west);
  \draw[flowarrow] (update.north) -- (formula.south);
  \draw[flowarrow] (formula.north) -- (inc.south);
  \draw[flowarrow] (inc.west) -- (loopjoin);

  \draw[flowarrow] (test.south) -- node[right=3pt] {否} (output.north);
  \draw[flowarrow] (output.south) -- (finish.north);

  \path[use as bounding box] (-2.45,.45) rectangle (6.15,12.20);
\end{tikzpicture}
